%%%%%%%%%%%%%%%%%%%%%%%%%
\spellentry{Quench}
%%%%%%%%%%%%%%%%%%%%%%%%%

\tagline{Transmutation}
\begin{description*}
\item[Level:] Drd 3
\item[Casting Time:] 1 standard action
\item[Range:] Medium (100 ft. + 10 ft./level)
\item[Duration:] Instantaneous
\item[Saving Throw:] None or Will negates (object)
\item[Spell Resistance:] No or Yes (object)
\item[Components:] V, S, DF
\end{description*}

\textit{Quench} is often used to put out forest fires and other conflagrations.
It extinguishes all nonmagical fires in its area. The spell also dispels any fire
spells in its area, though you must succeed on a dispel check (1d20 +1 per caster
level, maximum +15) against each spell to dispel it. The DC to dispel such spells
is 11 + the caster level of the fire spell.
Each elemental (fire) creature within the area of a \textit{quench} spell takes
1d6 points of damage per caster level (maximum 15d6, no save allowed).
Alternatively, you can target the spell on a single magic item that creates or
controls flame. The item loses all its fire-based magical abilities for
1d4 hours unless it succeeds on a Will save. (Artifacts are immune to this effect.)